\scanpage{4}
En d'autres termes, on a $f\overset{\mathrm{isot.}}{\sim}g$ ssi il existe un rev. \inserted{fini} étale (unique à isom. unique près) $Y'$ de $Y$, de telle façon que l'on ait
\[
(6)\qquad f=pf_0,\qquad g=pg_0
\]
\margin{$p:Y'\mathbin{-}Y$ la projection}
où
\[
(7)\qquad f_0:X\xrightarrow{\sim}Y'\qquad g_0:X\simeq Y'
\]
sont des homéomorphismes de $X$ \unsure{sur|en} $Y'$ qui sont isotopes, ou ce qui revient au $=$, tels que l'on ait
\[
(8)\qquad g_0=f_0\circ u
\]
où
\[
(9)\qquad u:X\xrightarrow{\sim}X
\]
\unsure{où} est un homéom. de $X$ sur lui $=$ isot. à l'identité. On aura donc, par (6) et (8),
\[
g=pg_0=p(f_0\circ u)=(pf_0)u=f\circ u
\]
Donc on trouve

\uline{Proposition 1} Pour que deux \deleted{applications} \inserted{morphismes} fini\deleted{\unsure{e}}s étales $f,g:X\longrightarrow Y$ soient isotopes, il f. et s. qu'ils soient \unsure{congrus|conjugués} sous le groupe des autom. de $X$ isotopes à l'identité; i.e $g=fu$, où $u$ autom. de $X$ isotope à l'identité.

Soit un \deleted{\unsure{donné} d'} morphisme $f:X\longrightarrow Y$ dans la catégorie isotopique (1), peut-on le
\margin{3}
